% !TeX root = ../泛函分析笔记.tex
% 由主文档通过 \input 载入；请编译主文档。
\setcounter{chapter}{0} % 保留原书章号。
\chapter{绪论：从线性代数到泛函分析}
\label{chap:fa-introduction}

\section{研究对象与基本问题}

\emph{泛函分析（functional analysis）}研究带有拓扑、范数或内积结构的线性空间，以及这些空间上的函数与算子。它将有限维线性代数中的向量、矩阵和线性方程推广到函数空间，使微分方程、积分方程、函数逼近与变分问题能够在共同的框架中讨论。

全书以三个问题组织内容：首先，如何选择容纳未知函数的空间，并保证极限仍留在该空间中；其次，怎样判断一个方程是否有解、解是否唯一，以及解对数据是否连续依赖；最后，当有限维的紧性与对角化方法失效时，如何借助弱收敛、紧算子和谱理论恢复可用的结构。

学习这些问题时，可以把函数暂时看成有无穷多个坐标的向量：数列的坐标用自然数编号，函数的取值用自变量编号。但这种类比只解释线性运算，不能替代分析。有限维中逐坐标收敛就是向量收敛；无限维中，每个坐标都变小，全部坐标合起来的误差却可能不变。因而本书始终同时追问两个问题：对象之间有什么代数关系，以及所用极限究竟在哪一种意义下成立。

\section{从一个积分方程开始}
考虑未知连续函数 $u:[0,1]\to\mathbb R$ 满足
\begin{equation}
 u(t)=t+\frac12\int_0^t u(s)\,\mathrm ds.
 \label{eq:fa-intro-integral}
\end{equation}
如果已经知道 $u$ 连续，右端可微，就能求导得到 $u'=1+u/2$、$u(0)=0$，从而猜到 $u(t)=2(e^{t/2}-1)$。不过，更一般的积分方程未必能这样求导，微分后的方程也未必有显式解。泛函分析提供的办法是直接研究整个函数所组成的空间。

从 $u_0(t)=0$ 出发，把已知函数代入右端，依次得到
\[
 u_1(t)=t,\qquad u_2(t)=t+\frac{t^2}{4},\qquad
 u_3(t)=t+\frac{t^2}{4}+\frac{t^3}{24}.
\]
记右端变换为 $F$，用 $\|u\|_\infty=\max_{0\leq t\leq1}|u(t)|$ 测量函数大小，则
\[
 \|F(u)-F(v)\|_\infty
 \leq\frac12\|u-v\|_\infty.
\]
也就是说，每次代入至少把两条输入曲线的最大差缩小一半。因此相邻迭代的误差形成可求和的几何级数，$(u_n)$ 是 Cauchy 序列。连续函数空间在这一距离下完备，保证极限仍为连续函数；$F$ 连续，保证极限满足原方程；同一个收缩估计还保证两个解只能相同。这正是第\ref{chap:fa-metric}章压缩映射原理的具体模型。

这个例子中，空间回答“解住在哪里”，范数回答“怎样测量误差”，算子回答“方程怎样变换函数”，完备性回答“近似过程会不会跑出空间”。后面处理最佳逼近时，用投影代替迭代；处理能量极小化时，用弱收敛寻找极限；处理积分方程时，用谱判断哪些参数会破坏可解性。方法虽然不同，选择空间、建立估计和传递极限的次序保持不变。

\section{课程层次与先修知识}

本书面向学过数学分析与高等代数的读者，主体定位为本科高年级课程。先修知识包括极限、连续性、函数列、级数、线性空间、线性映射、内积与有限维特征值理论。涉及可测函数空间时，需要补充测度与 Lebesgue 积分；涉及一般谱理论时，需要复分析中的解析函数与 Liouville 定理。

\begin{description}
  \item[入门基础] 第\ref{chap:fa-metric}章至第\ref{chap:fa-baire}章建立度量空间、Banach 空间、Hilbert 空间及基本定理，是第一次学习的重点。
  \item[核心理论] 第\ref{chap:fa-weak}章至第\ref{chap:fa-compact-spectral}章讨论弱拓扑、弱紧性、算子谱与紧算子，构成完整基础课程的后半部分。
  \item[应用专题] 第\ref{chap:fa-variational}、\ref{chap:fa-fixed-point}章将抽象理论用于变分问题、弱解和非线性方程，可按课程需要选读。
  \item[进阶专题] 第\ref{chap:fa-unbounded}章至第\ref{chap:fa-semigroups}章介绍无界算子、一般谱定理和算子半群，作为进一步学习的入口。
\end{description}

附录\ref{app:fa-measure}补充测度与积分，附录\ref{app:fa-topology}补充一般拓扑与集合论，附录\ref{app:fa-complex}补充复分析工具，附录\ref{app:fa-solutions}提供部分习题解答。尚未学习实分析的读者，可先用连续函数空间与数列空间理解主体概念，再结合附录进入可测函数空间。

\section{全书章节总览}

以下顺序兼顾知识依赖与学习动机。首次学习以基础与核心部分为主，应用部分择要阅读，进阶部分用于了解后续方向。

\begin{longtable}{L{0.08\textwidth}L{0.32\textwidth}L{0.48\textwidth}}
\toprule
章节 & 主题 & 核心问题 \\
\midrule
\endhead
\ref{chap:fa-introduction} & 从线性代数到泛函分析 & 从有限维到无限维，哪些方法仍然有效？ \\
\ref{chap:fa-metric} & 度量空间、完备性与紧性 & 极限何时存在，序列何时能够提取收敛子列？ \\
\ref{chap:fa-banach} & 赋范线性空间与 Banach 空间 & 怎样选择函数空间，怎样识别无限维现象？ \\
\ref{chap:fa-operators} & 有界线性算子与对偶空间 & 如何用范数控制线性映射，如何用标量观测向量？ \\
\ref{chap:fa-hilbert} & Hilbert 空间与正交展开 & 投影、最佳逼近和正交坐标如何推广？ \\
\ref{chap:fa-hahn-banach} & Hahn--Banach 定理与凸集分离 & 局部定义的泛函如何延拓，如何分离点与凸集？ \\
\ref{chap:fa-baire} & Baire 纲定理与算子基本定理 & 完备性怎样产生一致估计和逆映射的连续性？ \\
\ref{chap:fa-weak} & 弱拓扑、弱星拓扑与收敛 & 放宽收敛要求以后，哪些线性结构仍被保留？ \\
\ref{chap:fa-reflexive} & 弱紧性与自反空间 & 何时能从有界序列中提取弱收敛子列？ \\
\ref{chap:fa-spectrum} & 有界算子的谱理论 & 无限维中的不可逆性为何不限于特征值？ \\
\ref{chap:fa-compact} & 紧算子与 Fredholm 理论 & 哪些无限维方程仍具有有限维式的可解性结构？ \\
\ref{chap:fa-compact-spectral} & 紧自伴算子的谱分解 & 哪些算子能按正交特征向量展开？ \\
\ref{chap:fa-variational} & 变分方法与方程的弱解 & 如何由能量、弱紧性与强制性证明解的存在？ \\
\ref{chap:fa-fixed-point} & 不动点方法与非线性方程 & 非线性方程何时可由压缩性或紧性求解？ \\
\ref{chap:fa-unbounded} & 无界算子与定义域 & 微分算子的定义域为何是算子的一部分？ \\
\ref{chap:fa-general-spectral} & 一般谱定理与函数演算 & 无法用离散特征向量表示时，如何进行谱分解？ \\
\ref{chap:fa-semigroups} & 算子半群与演化方程 & 如何把矩阵指数推广为无限维时间演化？ \\
\bottomrule
\end{longtable}

\section{学习路线与贯穿全书的例子}

基本路线依次是空间结构、连续线性映射、基本存在性定理、弱收敛与紧性、算子谱及应用。Hilbert 空间的几何方法可以先于一般对偶理论学习；Hahn--Banach 定理与 Baire 纲定理各有独立的证明来源，不能将二者混为一条推导链。

全书反复使用连续函数空间、数列空间、可积函数空间，以及移位算子、对角算子、积分算子和微分算子。第一次出现时给出定义，此后在新的结构下重新考察同一对象。这样能够具体比较：同一线性映射换用范数后可能失去连续性；同一有界序列在弱拓扑下可能收敛；同一微分表达式采用不同边界条件会给出不同算子。

每章习题按概念辨析、证明训练、计算与应用组织。解题时先写清空间、标量域、范数或拓扑，再检查所用定理的假设。反例与正面结论同样重要，尤其要追踪完备性、闭性、有限维性、紧性和自反性分别用在何处。

第一次阅读时，可先以 $C([0,1])$、$\ell^1$、$\ell^2$ 和 $c_0$ 为四个主要模型，把涉及一般测度空间的段落留待第二遍阅读。弱紧性一章先掌握 Hilbert 空间中的子列构造，再读一般拓扑证明；谱理论一章先理解对角算子与乘法算子，再进入依赖复分析的谱半径证明。应用章只调用已经说明的结论；进阶章中引用的谱定理和生成定理均明确标出证明范围。

\section{记号与约定}

标量域记为 $\mathbb K$，取 $\mathbb R$ 或 $\mathbb C$。一般赋范线性空间记为 $X,Y$，Hilbert 空间记为 $H$，零向量记为 $0$。$\|x\|_X$ 表示 $X$ 中向量 $x$ 的范数；上下文明确时省略下标。$B_X$ 表示闭单位球 $\{x\in X:\|x\|_X\leq 1\}$。对集合 $M\subset X$，$\overline M$ 默认表示范数闭包，$\operatorname{span}M$ 表示有限线性组合的集合。

连续线性映射 $T:X\to Y$ 的全体记为 $\mathcal B(X,Y)$，并令 $\mathcal B(X)=\mathcal B(X,X)$；其核与值域分别记为 $\ker T$、$\operatorname{Ran}T$。连续线性泛函构成的对偶空间记为 $X^*$，双对偶记为 $X^{**}$。恒等算子记为 $I$。

复内积 $\langle x,y\rangle$ 约定对第一个变量线性、对第二个变量共轭线性。Banach 空间算子的对偶算子记为 $T':Y^*\to X^*$，Hilbert 空间算子的伴随记为 $T^*$，以免混淆两种构造。涉及弱星拓扑时使用 $\sigma(X^*,X)$，涉及弱拓扑时使用 $\sigma(X,X^*)$；完整定义见第\ref{chap:fa-weak}章。

\section{参考资源与本书的组织方式}

核心主题的覆盖范围参考 MIT 的泛函分析入门课程大纲\cite{mit-syllabus}，相关公开讲义可由课程阅读页\cite{mit-notes}获取。本书采用自己的章节顺序，将弱紧性与自反空间单列，并另设应用和进阶部分。章节总览用于组织知识，不等同于任何一门课程的考试范围。

